\documentclass[11pt,a4paper]{article}
\usepackage[margin=22mm]{geometry}
\usepackage[T1]{fontenc}
\usepackage{lmodern}
\usepackage{amsmath,amssymb,mathtools}
\usepackage{microtype,xcolor}
\definecolor{linkblue}{RGB}{20,57,102}
\usepackage[colorlinks=true,linkcolor=linkblue,citecolor=linkblue,urlcolor=linkblue]{hyperref}
\usepackage{bookmark}
\setlength{\parindent}{0pt}
\setlength{\parskip}{4pt plus 1pt}
\emergencystretch=2em
\newcommand{\R}{\mathbb R}
\newcommand{\C}{\mathbb C}
\newcommand{\tr}{\operatorname{tr}}
\newcommand{\EPG}{E_P^{\mathrm G}}
\newcommand{\ket}[1]{\lvert#1\rangle}
\newcommand{\bra}[1]{\langle#1\rvert}
\newcommand{\heading}[1]{\par\vspace{5pt}{\large\bfseries #1}\par\vspace{1pt}}
\hypersetup{pdftitle={Removing excess modes in Gaussian purification: an expert brief},pdfsubject={The transfer, equality and removal mechanism; companion to manuscript v0.1.1}}

\begin{document}
{\LARGE\bfseries Removing excess modes in Gaussian purification}\par
\vspace{4pt}
{\small Expert brief v0.1.1 \quad Companion to manuscript v0.1.1}

\heading{The result and its scope}
For a bipartite Gaussian state $\rho_{AB}$ with $n_A$ and $n_B$ modes, Gaussian entanglement of purification minimizes $S(AA')$ over pure Gaussian extensions and \emph{all finite} auxiliary mode counts. The accompanying manuscript~[1] proves
\begin{equation*}
\EPG(\rho_{AB})
=\min_{\substack{\psi_{ABA'B'}\ \text{pure Gaussian},\ \psi_{AB}=\rho_{AB}\\
n_{A'}=n_A,\ n_{B'}=n_B}}S(\psi_{AA'}).
\end{equation*}
Here bosonic states are normal with finite covariance; fermionic states are parity-invariant quasifree states. Pure factors, degenerate spectra and fermionic covariance zero modes are included. The main fermionic optimization allows both total-parity sectors.

This is the mode-count assertion of Windt, Jahn, Eisert and Hackl's minimum purification conjecture~[2, Sec.~6.1, Conjecture~2]. Their Section~6.3 explicitly states that the matched-mode restriction can be made without loss of generality and gives a canonical purification. Building on these results, [1] gives an explicit variational comparison with every finite auxiliary size and proves attainment. The result fixes the size of the Gaussian optimization; it neither settles optimality among non-Gaussian purifications nor guarantees that a local optimizer finds the global minimum.

\heading{The key mechanism: transfer, equality, removal}
The local ingredient is a compression lemma. If a Gaussian state on $XC$ has $a$ modes in $X$ and $c>a$ modes in $C$, an auxiliary-only Gaussian coordinate change selects a one-mode subsystem $P$, with $C=\widetilde C\oplus P$, such that
\begin{equation*}
S(X\widetilde C)\le S(XC),
\qquad
\text{equality}\ \Longleftrightarrow\ 
\rho_{XC}=\rho_{X\widetilde C}\otimes\ket\chi\bra\chi_P.
\end{equation*}
The equality claim applies to \emph{the mode selected by the construction}; it is the part that makes deletion possible. Fermionic factorization uses the graded tensor product.

Write $n=n_A+n_B$. For fixed total auxiliary count $M\ge n$, minimize over all extensions \emph{and all left/right assignments}; denote the attained minimum by $e_M$. Its existence is justified on the next page.

If $M>n$, at least one side has excess modes. Suppose it is $A'$. Apply the lemma to $AA'$ and first \emph{transfer} the selected mode:
\begin{equation*}
A\widetilde A'P\ \big|\ BB'
\quad\longrightarrow\quad
A\widetilde A'\ \big|\ B(PB').
\end{equation*}
The full state, physical marginal and auxiliary total are unchanged. Thus the new cut is an admissible competitor for the \emph{same} minimum $e_M$. Minimality gives the first inequality below, and compression gives the second:
\begin{equation*}
e_M\ \le\ S(A\widetilde A')\ \le\ S(AA')\ =\ e_M.
\end{equation*}
Both must be equalities. By the lemma, $P$ has a pure marginal, so it factorizes from the \emph{entire} purification. Only now can it be removed at no entropy cost.

Removal gives $e_{M-1}\le e_M$; padding by a pure mode gives the reverse inequality. Hence $e_M=e_{M-1}$ for $M>n$. Excess on the $B'$ side is treated using global purity, $S(AA')=S(BB')$.

At total $n$, successive transfers from the excess side reach the matched split without raising entropy. Any extension with fewer than $n$ auxiliaries can be padded to total $n$. This completes the comparison with \emph{every finite} auxiliary count.

{\small\textbf{AI contribution.} AI systems carried out problem selection, proof development and computational checks in an author-designed workflow. Subsequent AI reviews do not constitute human peer review. Details: [1].}

\newpage
\heading{Inside the compression lemma}
Let $d=a+c$ and let $\mathcal C\subset\R^{2d}$ be the $2c$-dimensional auxiliary coefficient space. The joint Gaussian normal form supplies a compatible complex structure $J$, with $J^2=-I$, even for a mixed state. The intersection $\mathcal C\cap J\mathcal C$ is $J$-invariant, and
\begin{equation*}
\dim_{\R}(\mathcal C\cap J\mathcal C)\ge 4c-2d=2(c-a)>0.
\end{equation*}
It therefore contains a complex line $P=\operatorname{span}_{\R}\{v,Jv\}$ lying wholly in the auxiliaries. Compatibility makes $P$ symplectic in the bosonic case. Its retained complement $U=X\oplus\widetilde C$ is $J$-invariant: use the orthogonal complement for fermions, the symplectic complement for bosons. The entropy on $U$ is obtained from a Hermitian compression. Only the auxiliary split is physically changed; the joint normal form is used to calculate its spectrum.

For fermions, $\Gamma=JK$ with $0\preceq K\preceq I$ and $[J,K]=0$. On the complex space defined by $J$, $K$ is Hermitian. Its compression $K_U$ gives the reduced covariance $\Gamma_U=J_UK_U$, where $J_U=J|_U$. If $\lambda_j$ and $\mu_j$ are the ascending eigenvalues of $K$ and $K_U$, interlacing gives $\lambda_j\le\mu_j\le\lambda_{j+1}$. The single-mode entropy $f$ is strictly decreasing on $[0,1]$, nonnegative, and zero only at $1$. Thus
\begin{equation*}
S(XC)-S(X\widetilde C)
=f(\lambda_d)+\sum_{j=1}^{d-1}\bigl[f(\lambda_j)-f(\mu_j)\bigr]\ge0.
\end{equation*}
Equality makes every term vanish: $\lambda_d=1$ and $\mu_j=\lambda_j$. For a complex unit vector $p$ spanning the removed line,
\begin{equation*}
\langle p,Kp\rangle=\tr_{\C}K-\tr_{\C}K_U=1,
\qquad I-K\succeq0\ \Longrightarrow\ Kp=p.
\end{equation*}
Consequently the selected mode is pure and its cross covariance vanishes. All eigenvalues and traces here count each mode once, in complex dimension.

For bosons, the corresponding Hermitian matrix is the Williamson diagonal $D\succeq I$ in its complex representation. Its eigenvalues $\nu_j$ interlace those of the compression as $\nu_j\le\mu_j\le\nu_{j+1}$. The single-mode entropy $b$ is strictly increasing and unbounded on $[1,\infty)$, with $b(1)=0$, so one uses the \emph{upper} bound $\mu_j\le\nu_{j+1}$. Equality forces $\nu_1=1$ and $\mu_j=\nu_{j+1}$. For a unit vector $p$ spanning the selected line in Williamson coordinates, the trace identity gives $\langle p,Dp\rangle=1$. Positivity of $D-I$ gives $Dp=p$, again a decoupled pure mode. Full covariance details: [1, Secs.~3.1--3.2].

\heading{Why the fixed-total minimum exists}
Bosonic first moments can be removed by local displacements without changing entropy. All centered pure Gaussian extensions at fixed $M$ form the auxiliary orthogonal or symplectic orbit of a reference extension [1, App.~A]. Fixing the left auxiliary count $k$ thus reduces the problem to cuts $A\oplus W$, where $W$ is a real $2k$-plane in the auxiliaries. Fermionic cuts form a compact Grassmannian, and entropy is continuous.

Bosonic cuts require a nondegenerate restricted symplectic form. In an orthonormal frame of $A\oplus W$, let $V_R$ and $\Omega_R$ be the restricted covariance and commutator forms. The fixed reference covariance gives $V_R\succeq v_*I$ with $v_*>0$, and [1, Sec.~4] proves
\begin{equation*}
\nu_{\max}(R)\ge\frac{v_*}{\sigma_{\min}(\Omega_R)},
\qquad S(AW)\ge b\bigl(\nu_{\max}(R)\bigr)\longrightarrow\infty
\end{equation*}
at a degenerate cut. Here $\sigma_{\min}$ is the smallest singular value and $\nu_{\max}$ the largest symplectic eigenvalue. Finite-entropy sublevels avoid that boundary, so compactness yields a minimizer. Minimizing over the $M+1$ choices of $k$ gives $e_M$. The bound is needed separately at each finite $M$.

\vspace{3pt}
{\small
\textbf{Full argument and sources.}
[1] \emph{A Mode Bound for Gaussian Entanglement of Purification}, accompanying manuscript v0.1.1 (28 September 2026). Compression: Sec.~3, pp.~3--5; auxiliary orbit: App.~A, pp.~8--9; attainment: Sec.~4, pp.~6--7; reduction: Sec.~5, pp.~7--8; parity: App.~B, p.~10.

[2] B.~Windt, A.~Jahn, J.~Eisert and L.~Hackl, \emph{Local optimization on pure Gaussian state manifolds}, SciPost Phys.~\textbf{10}, 066 (2021), \href{https://doi.org/10.21468/SciPostPhys.10.3.066}{doi:10.21468/SciPostPhys.10.3.066}.

}
\end{document}
